\documentclass{article}
\usepackage[T1]{fontenc}
\usepackage{lmodern}
\usepackage{graphicx}
\usepackage{amsmath,amssymb}
\usepackage[a4paper,margin=2cm]{geometry}
\usepackage[absolute,overlay]{textpos}
\setlength{\TPHorizModule}{1mm}
\setlength{\TPVertModule}{1mm}
\usepackage[indonesian]{babel}
\usepackage{hyperref}
\setlength{\emergencystretch}{2em}
% unit: Halaman 13

\begin{document}

% === Halaman 13 (python/native) ===

Bukti bahwa pesan m dapat diungkap kembali dari pasangan cipherteks a dan b: 

Dari persamaan (2): a = gk mod p   →  gk a (mod p)


Pangkatkan kedua ruas dengan x : gxk ax (mod p)


Dari persamaan (3):  b = ykm mod p   → ykm  b (mod p) 


Dari persamaan (1):  y = gx mod p  →  gx  y (mod p), 

maka 

b/ax ykm/ax (mod p)

gxkm/gxk (mod p)

m (mod p)

yang berarti bahwa plainteks m dapat diungkap kembali dari cipherteks a dan b.

13

\end{document}
